%-*-tex-*-
\ifundefined{writestatus} \input status \relax \fi %
\chcode{graph}
\def\cqu{}


\chapterhead{graph}{\TeX Graph}



The graphics package that is a part of \intex\ 
adds the following
capabilities 
\bl
\li $\bullet$ The production of a complete paper/report requiring no
additional insertions (except perhaps colour photographs if required).

\li $\bullet$ The ability to insert a file, for example a plot or halftone,
in Encapsulated PostScript Format (EPSF), and to automatically leave enough
space in the text. 

\li $\bullet$ The mixing of a vector graphics and \TeX\
text. In fact \TeX\ text may now be printed at any arbitrary angle.

\el

The package requires a sensible PostScript driver. The basic requirements
for the driver are as follows:
\bl
\lb The |\special| commands and the other information in the .dvi file are
processed in the order they are found in the .dvi file. This is necessary
in order to be able to rotate tables and arbitrary text \dots\ among other
things.
\lb The driver does not perform PostScript {\tt save} and {\tt restore}
commands around |\special| commands.
\lb The driver uses the code {\tt ps:} is used to signal a file name and
the code {\tt ps::} to signal inline code.
\lb The driver sets the {\tt currentpoint} before entering a |\special|.
\el 

\shead{gcomlist}{Command List - \TeX Graph}
\begintwocolumn
\ext|\absoluteposition|
\ext|\absolutescale|
\ext|\arcc|
\ext|\arcn|
\ext|\arrowheadscale|
\ext|\arrowheadsize|
\ext|\arrowheadtype|
\ext|\avec|
\ext|\beginsegment|
\ext|\beginTeXgraph \btg|
\ext|\cavec|
\ext|\cav|
\ext|\centergraph|
\ext|\centerplotfile|
\ext|\clvec|
\ext|\clv|
\ext|\endsegment|
\ext|\endTeXgraph \etg|
\ext|\graphdim|
\ext|\hgraphsize|
\ext|\hgraphskip|
\ext|\hpix|
\ext|\htext|
\ext|\includefile|
\ext|\incscfile|
\ext|\incscmvfile|
\ext|\larc|
\ext|\lcir|
\ext|\linecap|
\ext|\linejoin|
\ext|\lpatt|
\ext|\lvec|
\ext|\move|
\ext|\penwidth|
\ext|\relativeposition|
\ext|\relativescale|
\ext|\rstext|
\ext|\rtext|
\ext|\segmentscale|
\ext|\unitscale|
\ext|\vgraphsize|
\ext|\vgraphskip|
\ext|\vpix|
\ext|\vtext|
\endtwocolumn

\shead{plotfile}{The Inclusion of Encapsulated PostScript}
All the commands that are used to include encapsulated PostScript files
read the {\tt Boundingbox} and leave enough space in the text for the
graphics. The {\tt Boundingbox} may be at either the beginning or the end
of the  file. However, if the {\tt Boundingbox} is at the the end, it might
take \TeX\ some time to read the entire file to find it.

The simplest command for including a graphics file is 
$$ 
|\centerplotfile{<file name>}|.
$$
For instance the (old) IEEE Communications Society
logo was brought in simply by saying 
\beginttverbatim \ninepoint \tt
\1medskip
\1hrule
\1centerplotfile{ieeelogo.ps}
\1hrule
\1medskip
\endttverbatim
right here.
\medskip
\hrule 
\centerplotfile{ieeelogo.ps}
\hrule
\medskip
The |\hrule| commands were put in to show how much space was left by \TeX
Graph. 
A slightly more complicated example illustrating both rotation and scaling
is given below.
\bigskip
\hrule
\centergraph{ \btg
              \incscfile f:{ieeelogo.ps} sc:.3 d:45 
              \etg}
\hrule
\bigskip
The code for this is 
\beginttverbatim
\1hrule
\1centergraph{ \1btg
              \1incscfile f:{ieeelogo.ps} sc:.3 d:45 
              \1etg}
\1hrule
\endttverbatim

Notice that the computation of the space to leave was incorrect. This is
due to the difficulty in performing the computations for obtaining the
bounding box for an arbitrarily rotated and scaled figure. 

The most general
inclusion command is |\incscmvfile| which allows for arbitrary scaling,
rotation, and offset. Finally, the rotation, scaling, and placement may be
applied around any of the nine text box reference points described in
Section~\ref{tgraph}.


\shead{tforms}{\TeX Graph Basics}
\TeX Graph has a number of capabilities that range from primitive to
powerful. The following list of classes is for reference only. More detail
will follow with examples in the various sections to follow. This section may
be omitted on first reading. 
\dssshead{\TeX Graph Classes}
\bl
\lb Drawing primitives consist of commands for circles,
circular arcs, line vectors, vectors with arrowheads, spline curves, spline
curves with arrowheads, line dotting commands, and various shades of gray
for fills.

\lb There are commands for changing whether the distances specified
are to be interpreted absolutely, with respect to an origin or relatively
with respect to the present position. 

\lb There is a command to change the (pen)width of the drawn lines
and to change the interpreted scale of the units of distance. 

\lb There is a command to change the default units. However, since
arbitrary scaling of units is allowed, this command is of limited use except
for those that visualize distances in different dimensions. The dimensions
available are those available in \intex.

\lb There is  a ``graphical segment''  which can be
placed at any specific location. Distances inside a graphical segment are
interpreted with respect to the segment origin. The distances inside a
segment may be either relative or absolute with respect to the segment
origin. There is a {\it segment scale} which allows for segments to be scaled
relatively with respect to an overall graph. 

\lb Text, with all the capabilities of \intex\ may be written
at any arbitrary angle.

\lb \TeX Graph keeps track of the maximum extent of a graph. These
are available after the graph is terminated in the variables |\hgraphsize|
and |\vgraphsize|. |\hgraphskip| and |\vgraphskip| will cause a horizontal or
vertical skip of the appropriate graphsize. {\bf Unfortunately, these values
may be incorrect for a number of reasons. } These are
\bsl
\li $\circ$ Text may stick out the edge of the graph.
\li $\circ$ A circle or arc center is outside the main graph area.
\li $\circ$ The visible part of a circle or arc may be too far to the right.
\li $\circ$ A vertical segment may have inadvertently reset the maximum
positions.
\esl
In some cases it may be necessary to make a fake move to the lower right hand
corner of a graph to tell \TeX Graph the correct size \dots\ after you have
seen a preliminary version and know what it should be.

\lb All of \TeX's definitional capabilities are available. These
can be used to replace a complex structure by a single command or
to change the surface syntax of \TeX Graph. 

\lb Explicit use of grouping in \TeX Graph, through $\{ \}$, or
|\begingroup|, |\endgroup| pairs is not allowed. This will foul up the
positional record keeping inside. Groups may be produced only through the
use of segments, namely the |\beginsegment|, |\endsegment| pair. 
\el


\shead{tgraph}{Simple \TeX Graphics}
\TeX Graph will allow for rather sophisticated graphics. You can use the 
definition capabilities of  \intex\ to create libraries of symbols 
for future use. 

Directions in \TeX Graph are described as horizontal, {\tt h:}, increasing
from left to right across the page and vertical, {\tt v:}, increaseing from
top to bottom {\bf down} the page. Thus the fundamental directions are 
$$
\centergraph{\btg
             \move h:.3 v:.3 
             \avec h:3 v:.3
             \move h:.3 v:.3 
             \avec h:.3 v:2 
             \textref h:C v:C
             \move h:1.5 v:.15 \htext{Horizontal -- {\tt h:<number>}}
             \move h:.15 v:1 \vtext{Vertical -- {\tt v:<number>}}
             \etg}
$$
\TeX Graph capabilities will be introduced slowly. 
The following diamond box 
$$
\centergraph{
   \btg 
     \move h:0 v:.5  % moves penup to (0,.5)
     \lvec h:.5 v:0  % draws a line to (.5,0)
     \lvec h:1 v:.5
     \lvec h:.5 v:1
     \lvec h:0 v:.5     
     \textref h:C v:C  % sets the reference point in center of text.
     \move h:.5 v:.5
     \htext {$\sum_{i=1}^n \rho_n$} % horizontal text box
    \etg}
$$
was put in place with
\beginttverbatim \ninepoint \tt
$$  % a convenient way to use displaymath spacing
\1centergraph{
   \1beginTeXgraphics % also known as \1btg
     \1move h:0 v:.5  % moves penup to (0,.5)
     \1lvec h:.5 v:0  % draws a line to (.5,0)
     \1lvec h:1 v:.5
     \1lvec h:.5 v:.1
     \1lvec h:0 v:.5     
     \1textref h:C v:C  % sets the reference point in the center.
     \1move h:.5 v:.5
     \1htext {$\1sum_{i=1}^n \1rho_n$} % horizontal text box
    \1endTeXgraphics  % also known as \1etg
     }
$$
\endttverbatim

The command |\centergraph{<...>}| uses the information about plot size to
center the graphics horizontally and to leave enough room vertically. 

Arguments for all \TeX Graph primitive commands  are delimited by spaces.
The units of the dimensions are implicit, but changable. The default is
inches. Thus {\tt h:<number>} is the distance in the horizontal (to right)
and {\tt v:<number>} is the distance in the vertical (down) direction. The
locations are absolute distances with respect to an origin in the upper left
hand corner. 

\shead{vecforms}{Simple Vector Forms}
This introduces some baisc drawing commands. There {\bf must} be 
either a space or an
end of line after each {\tt <coordinate>}\]. The dimensional units are
implicit. We usually use {\tt <..>} instead of the less economical {\tt
<coordinate>}.
\beginblockmode
\ext\@|\move h:<..> v:<..>|
\nbr
This is a move from the current position to the coordinates given by
{\tt h:<..> v:<..>}. The pen is up. It may be used to move the current
reference point or to describe a boundary when creating a filled polygon.

\mbr\ext\@|\lvec h:<..> v:<..>|
\nbr
This draws a line from the current position to {\tt h:<..> v:<..>}.

\mbr\ext\@|\clvec h1:<..> v1:<..> h2:<..> v2:<..> h3:<..> v3:<..>|
\nbr
This draws a Bezier spline curve from the current point to {\tt h3:<..>
v3:<..>} with respect to the control points {\tt h1:<..> v1:<..>} and 
{\tt h2:<..> v2:<..>}. These curves are very versatile and can create many
effects, including loops,  by varying the control points. 

\mbr\ext\@|\clv  (<h1> <v1>)  (<h2> <v2>)  (<h3> <v3>)|
\nbr
This command is identical to |\clvec| but with a less verbose syntax. 
\endblockmode

Lines may be either solid of various widths, or dotted with custom
patterns. In addition the ends of the lines may be of different shapes and
joined with different defaults. In most cases you will not have to change
the defaults in \TeX Graph but the capability is there. 

\beginblockmode
\ext\@|\penwidth <size>|
\nbr
This sets the width of the pen or line in the current units. The default is
is .015 (in). 

\mbr\ext\@|\linecap <integer>|
\nbr
This sets the end shapes of the lines using PostScript notation. The
default is 1 or a round cap. A value of 0 is a butt cap with the line cut
off square exactly to its length. A value of 2 is a square cap with the
line extended by a square cap the thikness of half its width. For narrow
lines the difference is negligible. 

\mbr\ext\@|\linejoin <integer>|
\nbr
This sets the form in which lines join (PostScript notation). The default
is 1 which corresponds to a round line join. A value of 0 is a mitre or
pointed join while a value of 2 is a beveled or flattened join. For thin
lines the differences are negligible.

\mbr\ext\@|\lpatt p:<pattern>|
\nbr 
This sets the dotting pattern of the line in PostScript notation. The
default is a solid line described by |{[]}|. The simplest form is |{[a]}|
where the line is black for {\tt a} pixels and white for 1. |{[a b]}| would
mean {\tt a} pixels black followed by {\tt b} pixels white. A specification
of |[a1 b1 a2 b2]| would be a cyclic pattern with the line black for the
{\tt a?} pixels and white for the corresponding {\tt b?}.

\endblockmode

\TeX Graph alos supplies lines with arrowheads. These may be combined
dotting patterns given above but a decision will have to be made as to
whether you want the arrowhead itself dotted. 
There are
three styles of arrowheads 
         \btg \move h:0 v:-.05 \avec h:.3 v:-.05
           \arrowheadtype t:O 
           \avec h:.6 v:-.05
           \arrowheadtype t:F
           \avec h:.9 v:-.05
         \etg \hgraphskip\ 
triangle {\tt T}, open {\tt O},
and filled {\tt F}. The size may also be varied. There are no restriction on
the direction  of the arrows. The commands are

\beginblockmode
\ext\@|\avec h:<..> v:<..>|
\nbr
This draws a vector with an arrow head from the current position to
{\tt h:<..> v:<..>}.

\mbr\ext\@|\cavec h1:<..> v1:<..> h2:<..> v2:<..> h3:<..> v3:<..>|
\nbr\ext\@|\cav (<h1> <v1>) (<h2> <v2>) (<h3> <v3>) |
\nbr
These are the spline versions of the vectors with arrowheads. 

\mbr\ext\@|\arrowheadtype t:<type>|
\nbr
The  {\tt <type>} is {\tt T O F} for triangle, open, or filled respectively.

\mbr\ext\@|\arrowheadsize l:<length> w:<width>|
\nbr
The {\tt<length>} and {\tt<width>} are in the current units. The default
is 0.16in by 0.04in. This size is {\tt independent} of the current scales.
This prevents arrowheads from disappearing as a graph is scaled.

\mbr\ext\@|\arrowheadscale|
This makes the arrowheadsize sensitive to the segment and unit scales in
force when it is called. It is not global but is affected by the grouping
imposed by |\beginsegment| and |\endsegment| pairs.

\endblockmode

For most applications,  the |\arrowheadsize| and the |\arrowheadscale|
will not be needed. 

The following example of a (partial) block diagram uses these commands. 
$$
\centergraph{\btg
        \move h:0 v:.3 
        \avec h:.5 v:.3
        \move h:.5 v:.6 % moves to lower lefthand corner of box
        \lvec h:.5 v:0  % four sides 
        \lvec h:1.5 v:0
        \lvec h:1.5 v:.6        
        \lvec h:.5 v:.6
          \etg}
$$
This was produced by the following:
\beginttverbatim \ninepoint \tt
\1centergraph{\1btg
        \1move h:0 v:.3 
        \1avec h:.5 v:.3 % arrow vector leading into box
        \1move h:.5 v:.6 % moves to lower lefthand corner of box
        \1lvec h:.5 v:0  % four sides 
        \1lvec h:1.5 v:0
        \1lvec h:1.5 v:.6        
        \1lvec h:.5 v:.6
          \1etg}
\endttverbatim

We can change the type of arrowhead on this and all subsequent arrow vectors
by adding adding |\arrowheadtype t:O| before the |\avec ...| to give
$$
\centergraph{\btg
        \move h:0 v:.3 
        \arrowheadtype t:F % changes default filled
        \avec h:.5 v:.3 % arrow vector leading into box
        \move h:.5 v:.6 % moves to lower lefthand corner of box
        \lvec h:.5 v:0  % four sides 
        \lvec h:1.5 v:0
        \lvec h:1.5 v:.6        
        \lvec h:.5 v:.6
          \etg}
$$
This was produced by 
\beginttverbatim \ninepoint \tt
\1centergraph{\1btg
        \1move h:0 v:.3 
        \1arrowheadtype t:F % changes default filled
        \1avec h:.5 v:.3 % arrow vector leading into box
        \1move h:.5 v:.6 % moves to lower lefthand corner of box
        \1lvec h:.5 v:0  % four sides 
        \1lvec h:1.5 v:0
        \1lvec h:1.5 v:.6        
        \1lvec h:.5 v:.6
          \1etg}
\endttverbatim

\beginblockmode
\ext\@|\unitscale <scale>|
\nbr 
This command causes all subsequent units in this or enclosed segments to be
scaled by the {\tt <scale>}. Thus |\unitscale 2.54| would have the effect
of changing a figure where the units were originally in inches to be in
centimetres. 
\endblockmode


Thus the entire graph is reduced to half size by adding |\unitscale .5|\] at
the very beginning.
$$
\centergraph{\btg
        \unitscale .5
        \move h:0 v:.3 
        \arrowheadtype t:F % changes default filled
        \avec h:.5 v:.3 % arrow vector leading into box
        \move h:.5 v:.6 % moves to lower lefthand corner of box
        \lvec h:.5 v:0  % four sides 
        \lvec h:1.5 v:0
        \lvec h:1.5 v:.6        
        \lvec h:.5 v:.6
          \etg}
$$
Note that the arrowhead has not changed in size. The commands to do this are
\beginttverbatim \ninepoint \tt
\1centergraph{\1btg
        \1unitscale .5  % causes all subsequent units to be multiplied
                       % by .5 
        \1move h:0 v:.3 
        \1arrowheadtype t:F % changes default filled
        \1avec h:.5 v:.3 % arrow vector leading into box
        \1move h:.5 v:.6 % moves to lower lefthand corner of box
        \1lvec h:.5 v:0  % four sides 
        \1lvec h:1.5 v:0
        \1lvec h:1.5 v:.6        
        \1lvec h:.5 v:.6
          \1etg}
\endttverbatim
If the arrowhead size needed to be reduced in the same scale it is necessary
to add |\arrowheadscale| after the |\unitscale| and before the first
arrowvector. Thus we get
$$
\centergraph{\btg
        \unitscale .5  % causes all subsequent units to be multiplied
                       % by .5 
        \move h:0 v:.3 
        \arrowheadscale % makes arrowhead size sensitive to scalings. 
        \arrowheadtype t:F % changes default filled
        \avec h:.5 v:.3 % arrow vector leading into box
        \move h:.5 v:.6 % moves to lower lefthand corner of box
        \lvec h:.5 v:0  % four sides 
        \lvec h:1.5 v:0
        \lvec h:1.5 v:.6        
        \lvec h:.5 v:.6
          \etg}
$$
This is produced by 
\beginttverbatim \ninepoint \tt
\1centergraph{\1btg
        \1unitscale .5  % causes all subsequent units to be multiplied
                       % by .5 
        \1move h:0 v:.3 
        \1arrowheadscale % makes arrowhead size sensitive to scalings. 
        \1arrowheadtype t:F % changes default filled
        \1avec h:.5 v:.3 % arrow vector leading into box
        \1move h:.5 v:.6 % moves to lower lefthand corner of box
        \1lvec h:.5 v:0  % four sides 
        \1lvec h:1.5 v:0
        \1lvec h:1.5 v:.6        
        \1lvec h:.5 v:.6
          \1etg}
\endttverbatim

Finally we have an example with dotted Bezier splines. 
$$
\centergraph{ \btg \lpatt p:{[20 30 20]}
                   \lvec h:0 v:2 
                   \move h:0 v:0 
                   \lpatt p:{[20 20 10 10]}
                  \cav (1 1) (2 1) (3 0) 
                   \lpatt p:{[20 20]}
                  \move h:4 v:1
                  \cav (0 2) (5 2) (1 1)
               \etg
            }
$$

which was given by 
\beginttverbatim
\1centergraph{ \1btg \1lpatt p:{[20 30 20]}
                   \1lvec h:0 v:2 
                   \1move h:0 v:0 
                   \1lpatt p:{[20 20 10 10]}
                  \1cav (1 1) (2 1) (3 0) 
                   \1lpatt p:{[20 20]}
                  \1move h:4 v:1
                  \1cav (0 2) (5 2) (1 1)
               \1etg
            }
\endttverbatim
This is off center because the width computation includes the (5 2) control
point of the Bezier curve. 

\sshead{text}{Text in Graphs}
All text within \TeX Graph is actually typeset by \TeX. Since \TeX\ does
this so beautifully, it seems counter productive to do it otherwise. The
commands for including text at various angles are

\beginblockmode
\ext\@|\htext {<text>}| 
\nbr
This puts the text in a horizontal |\hbox| running from left to right.

\mbr\ext\@|\vtext {<text>}| 
\nbr
This puts the text in a vertical |\hbox| going upwards.

\mbr\ext\@|\rtext d:<degrees> t:{<text>}|
\nbr
This puts text in an |\hbox| with an arbitrary rotation of {\tt <degrees>}.

\mbr\ext\@|\rstext d:<degrees> sc:<scale> t:{<text>}|
\nbr
This puts text in an |\hbox| with an arbitrary rotation of {\tt <degrees>}
and scale {\tt <scale>}.

\endblockmode

Text is placed in an |\hbox|. This box is located at the point in the graph
where the |\?text| command is called. 
In order to make the exact placement of the
box in a graph easier, the box is actually oriented at any one of nine
reference points. The command is

\beginblockmode
\ext\@|\textref h:<L C R> v:<T C B>|
\nbr
The choices of {\tt L C R} are Left, Center, and Right respectively in the
horizontal and vertical positions. 
\endblockmode

The nine
choices are shown below:
\medskip
\def\bul{\htext {$\bullet$} }
\def\tref{\beginsegment
           \textref h:C v:C
           \lvec h:1 v:0 \bul 
           \lvec h:2 v:0 \bul 
           \lvec h:2 v:.5 \bul 
           \lvec h:2 v:1 \bul 
           \lvec h:1 v:1 \bul 
           \lvec h:0 v:1 \bul 
           \lvec h:0 v:.5 \bul 
           \lvec h:0 v:0 \bul 
           \move h:1 v:.5 \bul
          \endsegment}
\def\tlabels{\textref h:R v:C %left hand labels
             \htext{ h:L v:T } \lvec h:.2 v:0 \avec h:.6 v:.4 
             \move h:0 v:.9 \htext{ h:L v:C } \avec h:.6 v:.9 
             \move h:0 v:1.8 \htext{ h:L v:B } \lvec h:.2 v:1.8 \avec h:.6 v:1.4 
             \textref h:L v:C %right hand labels
             \move h:3.2 v:0 \htext{ h:R v:T } \lvec h:3 v:0 \avec h:2.6 v:.4 
             \move h:3.2 v:.9 \htext{ h:R v:C } \avec h:2.6 v:.9 
             \move h:3.2 v:1.8 \htext{ h:R v:B } \lvec h:3 v:1.8 \avec h:2.6 v:1.4 
             \textref h:C v:B % Center Top 
             \move h:1.6 v:0 \htext{ h:C v:T } \move h:1.6 v:.05 
                \avec h:1.6 v:.4 
             \textref h:L v:C % Center Center
             \move h:2.0 v:.7 \htext{ h:C v:C } \lvec h:1.85 v:.7 
                          \avec h:1.6 v:.9 
             \textref h:C v:T % Center Bottom
             \move h:1.6 v:1.8 \htext{ h:C v:B } \move h:1.6 v:1.75 
                           \avec h:1.6 v:1.4 
            }
\medskip
\centergraph{\btg \unitscale .7 
             \arrowheadsize l:.08 w:.02 \arrowheadtype t:O
             \tlabels
             \move h:.6 v:.4 \tref
             \etg}
\medskip
These reference points move with the box as it is rotated or scaled.

As an example we add the following $\rho=\sum_{i=1}^{n}\rho_i$ to the 
center of the block diagram used previously to give
$$
\centergraph{\btg
        \move h:1.0 v:.3 % moves to center of box
        \textref h:C v:C % sets text reference to dead center
        \htext {$\rho=\sum_{i=1}^{n}\rho_i$} % puts in text
        \move h:0 v:.3 
        \avec h:.5 v:.3 % arrow vector leading into box
        \move h:.5 v:.6 % moves to lower lefthand corner of box
        \lvec h:.5 v:0  % four sides 
        \lvec h:1.5 v:0
        \lvec h:1.5 v:.6        
        \lvec h:.5 v:.6
          \etg}
$$
The commands to do this are
\beginttverbatim \ninepoint \tt
\1centergraph{\1btg
        \1move h:1.0 v:.3 % moves to center of box
        \1textref h:C v:C % sets text reference to dead center
        \1htext {$\1rho=\1sum_{i=1}^{n}\1rho_i$} % puts in text
        \1move h:0 v:.3 
        \1avec h:.5 v:.3 % arrow vector leading into box
        \1move h:.5 v:.6 % moves to lower lefthand corner of box
        \1lvec h:.5 v:0  % four sides 
        \1lvec h:1.5 v:0
        \1lvec h:1.5 v:.6        
        \1lvec h:.5 v:.6
          \1etg}
\endttverbatim

Arbitrary rotation of text or graphics |\rtext|. 
This example uses |\rtext d:45 t:<text>| in place of the |\htext|.
$$
\centergraph{\btg
        \move h:1.0 v:.3 % moves to center of box
        \textref h:C v:C % sets text reference to dead center
        \rtext d:45 t:{$\rho=\sum_{i=1}^{n}\rho_i$} % puts in text
        \move h:0 v:.3 
        \avec h:.5 v:.3 % arrow vector leading into box
        \move h:.5 v:.6 % moves to lower lefthand corner of box
        \lvec h:.5 v:0  % four sides 
        \lvec h:1.5 v:0
        \lvec h:1.5 v:.6        
        \lvec h:.5 v:.6
          \etg}
$$
Note that the text spills out over the edges of the box. There is no
guarantee that the text will fit in the box. 
In fact if we  take the
horizontal example and add a |\unitscale .7 \arrowheadscale| the following
results:
$$
\centergraph{\btg
        \unitscale .7
        \arrowheadscale
        \move h:1.0 v:.3 % moves to center of box
        \textref h:C v:C % sets text reference to dead center
        \htext {$\rho=\sum_{i=1}^{n}\rho_i$} % puts in text
        \move h:0 v:.3 
        \avec h:.5 v:.3 % arrow vector leading into box
        \move h:.5 v:.6 % moves to lower lefthand corner of box
        \lvec h:.5 v:0  % four sides 
        \lvec h:1.5 v:0
        \lvec h:1.5 v:.6        
        \lvec h:.5 v:.6
          \etg}
$$

Note that the font is not scaled. If we use |\rstext d:0 sc:.8 t:<...>| we
get
$$
\centergraph{\btg
        \unitscale .7  
        \arrowheadscale
        \move h:1.0 v:.3 % moves to center of box
        \textref h:C v:C % sets text reference to dead center
        \rstext d:0 sc:.8 t:{$\rho=\sum_{i=1}^{n}\rho_i$} 
        \move h:0 v:.3 
        \avec h:.5 v:.3 % arrow vector leading into box
        \move h:.5 v:.6 % moves to lower lefthand corner of box
        \lvec h:.5 v:0  % four sides 
        \lvec h:1.5 v:0
        \lvec h:1.5 v:.6        
        \lvec h:.5 v:.6
          \etg}
$$


\shead{vertgraph}{Rotated Graphics and Text}
It is possible to place graphics or text at any arbitrary direction 
on a page. It is only
necessary to include it in a |\vtext| for vertical or an |\rtext| for an
arbitrary rotation. However, it should be remembered that
the reference point is that position on the page where the |\?text| was
entered {\bf and} that the text/graphics goes at the appropriate angle. 
The following shows a
simple table angled on the page. To do this correctly, requires that you
discover the size of the table. 

\def\sample{
   \btg 
   \lvec h:1 v:0 
   \incscfile f:ieeelogo.ps sc:.25 d:0
   \lvec h:0 v:.5
   \htext{\begintable
           \begintableformat
            &\center
           \endtableformat
           \-
           \br{\:|}\use{3} PROJECT --- Squidget Forms  \er{|}
           \-
           \br{\:|} Milestone | Due     |  Projected    \er{|}
           \-
           \br{\:|}    A      | Aug. 94 | Jan. 95       \er{|}
           \-
           \br{\:|}    B      | Dec. 94 | Dec. 95       \er{|}
           \-
           \br{\:|}    A      | Jan. 95 | ?             \er{|}
           \-
         \endtable} 
      \etg}

$$
\centergraph{
    \btg 
      \lvec h:0 v:1 % sets origin
      \rtext d:45 t:{\sample}
      \lvec h:2 v:2 % sets size of graph ... 
    \etg}
$$

This was done with the following set of commands. Note that a |\def| was used
to make the construction easier to follow. |\lvec| was used rather than
|\move| in order to make the origin movements visible.
\beginttverbatim \ninepoint \tt
\1def\1sample{
   \1btg 
   \1lvec h;1 v:0
   \1incscfile f:ieeelogo.ps sc:.25  d:0
   \1lvec h:0 v:.5
   \1htext{\1begintable
           \1begintableformat
            &\1center
           \1endtableformat
           \1-
           \1br{\1:|}\1use{3} PROJECT --- Squidget Forms  \1er{|}
           \1-
           \1br{\1:|} Milestone | Due     |  Projected    \1er{|}
           \1-
           \1br{\1:|}    A      | Aug. 94 | Jan. 95       \1er{|}
           \1-
           \1br{\1:|}    B      | Dec. 94 | Dec. 95       \1er{|}
           \1-
           \1br{\1:|}    A      | Jan. 95 | ?             \1er{|}
           \1-
         \1endtable} 
      \1etg}
\endttverbatim

\beginttverbatim \ninepoint \tt
\1medskip
\1centergraph{
    \1btg \1textref h:C v:C 
      \1lvec h:0 v:1 % sets origin
      \1rtext d:45 t:{\1sample}
      \1lvec h:2 v:2 % sets size of graph ... 
    \1etg}
\1medskip
\endttverbatim
In fact, this example was first
done as a separate file to learn the appropriate dimensions and then inputted
in the appropriate place. 

\sshead{arcs}{Circles, Arcs and Fills}
\TeX Graph uses the circle, arc, and fill commands of PostScript to
produce curves. PostScript uses a polygonal fill on the {\bf inside} of the
curve. This means that the direction of the arcs are important. 

The commands are 
\beginblockmode
\ext\@|\lcir r:<radius>|
\nbr
This is a circle with the center at the current loaction and a radius of
size {\tt <radius>} in the current dimensions.

\mbr\ext\@|\arcc r:<radius> sd:<degrees> ed:<degrees>|
\nbr
This is an ``unstroked''arc, ie equivalent to a |\move| in the counter
clockwise direction from {\tt sd:<..>} to {\tt ed:<..>}. If this is used in
a fill, the center of the circle is the beginning of the path. 

\mbr\ext\@|\arcn r:<radius> sd:<degrees> ed:<degrees>|
\nbr
This is the same as |\arcc| except that it is counter clockwise. 

\mbr\ext\@|\pfill p:<fraction>|
\nbr
This completes a path a fills the path with a shade of gray with {\tt p:0}
being black and {\tt p:1} being white. 
\endblockmode

The angles are in the directions shown here.
$$
\centergraph{
  \btg \unitscale .7
    \move h:1 v:2 
    \avec h:1 v:0
    \move h:0 v:1
    \avec h:2 v:1 
    \move h:1 v:2.05 \textref h:C v:T \htext{$90^\circ$}
    \move h:0 v:1 \textref h:R v:C \htext{$180^\circ$ }
    \move h:1 v:-.05 \textref h:C v:B \htext{$270^\circ$}
    \move h:2 v:1 \textref h:L v:C \htext{ $0^\circ,\ 360^\circ$ }    
   \etg}
$$

\def\fillbox #1 {\move h:.4 v:0
                  \beginsegment 
                   \absoluteposition
                   \lvec h:.5 v:0
                   \lvec h:.5 v:.5
                   \lvec h:0 v:.5
                   \pfill p:{#1}
                   \move h:.25 v:.55 \textref h:C v:T \htext{#1}
                 \endsegment
                 \move h:.4 v:0 % for symmetry
                 }

Some shades of gray are as follows:
$$
\centergraph{
\btg
\relativeposition
\fillbox 0
\fillbox .1
\fillbox .2
\fillbox .3
\fillbox .4
\move h:-4 v:.8
\fillbox .5
\fillbox .6
\fillbox .7
\fillbox .8
\fillbox .9
\etg}
$$

An example of a use of arcs and fill is given below. 
$$
\centergraph{\btg
              \arcc r:2 sd:0 ed:45  
              \pfill p:.5
              \arcc r:1 sd:0 ed:45 
              \pfill p:1  % fills the inside with white
              \move h:2 v:1.5 % to leave enough space
             \etg}
$$

This logo is given by 
\beginttverbatim \ninepoint \tt
\1centergraph{\1btg
              \1arcc r:2 sd:0 ed:45  
              \1pfill p:.5
              \1arcc r:1 sd:0 ed:45 
              \1pfill p:1  % fills the inside with white
              \1move h:2 v:1.5 % to leave enough space
             \1etg}
\endttverbatim
This shows that a fill of ``white'' (1) is opaque and obliterates the black
or gray underneath. Note the single line of dots from the origin. Apparently
there is some roundoff. This strategy is necessary in PostScript filling.


\shead{seg}{Advanced \TeX Graphics and Segments}
The most important advanced feature of \TeX Graph is the graphical segment. 
These are delimited by |\beginsegment| |\endsegment| pairs. They are used to
create graphical items that can be placed at any specific location in a
graph. This capability, combined with \TeX's ability to create definitions
allows for the creation of libraries of symbols. {\bf Segments are the only
groups that should be used within a} |\btg| |\etg| {\bf
pair.}\footnote{!}{This is almost true. Groups may be used freely within
the argument of either an |\??text|.} A simple example of the use
of this facility was the definition of a |\fillbox| to create the box with
the fill pattern shown in Section~\ref{arcs}. The definition was 

\beginttverbatim \ninepoint \tt
\1def\1fillbox #1 {\1move h:.4 v:0
                  \1beginsegment 
                   \1absoluteposition
                   \1lvec h:.5 v:0
                   \1lvec h:.5 v:.5
                   \1lvec h:0 v:.5
                   \1pfill p:{#1}
                   \1move h:.25 v:.55 \1textref h:C v:T \1htext{#1}
                 \1endsegment
                 \1move h:.4 v:0 % for symmetry
                 }
\endttverbatim

The relative position in outside the segment allows the boxes to be placed
beside each other with even spacing and no additional instructions. 

Each segment has a {\it segment reference point}. This is the location in
the enclosing segment or |\btg| |\etg|, which is really just a special
segment. {\bf Absolute positions} within a segment are actually relative to
the {\it segment reference point}. Relative positions are relative. 
The three
boxes below 
$$
 \centergraph{\btg \relativeposition
               \fillbox .3
               \fillbox .5
               \fillbox .7
              \etg}
$$
were produced by 
\beginttverbatim \ninepoint \tt
\1medskip
 \1centergraph{\1btg \1relativeposition
               \1fillbox .3
               \1fillbox .5
               \1fillbox .7
              \1etg}
\1medskip
\endttverbatim
Note that the space after the |#1| in the definition for |\fillbox| means
that the numbers are delimited by implicit space at the end of the line. If
this space had been left out, only the first digit in each of the pattern
numbers would have been read and the other would have been printed at some
odd place on the page. 

\sshead{segscale}{Segment Scaling} 
There are several scaling parameters that can be used with segments. In
addition to the |\unitscale| there is a |\segmentscale|. Roughly speaking the
|\segmentscale| is used to change the scale of the present and all enclosed
segments while the |\unitscale| is used to set the nominal size of segment
object. The major difference between the two is that segment scales
accumulate while unitscales do not. An example is obtained using the
|\fillbox| in a segment.  The following three boxes illustrate the effect.
$$
\centergraph{\btg \unitscale .3 % changes units 
               \segmentscale 2 % doubles scale of all enclosed segments
                           % actual scale is .3*2=.6
               \fillbox .5
               \move h:.8 v:0  % moves to right beyond box
               \beginsegment
                \unitscale .5 % now unit scale is .5
                \segmentscale .7 % reduces scale by .5 -- actual scale is
                              % .5*2*.7=.7 
                \fillbox .7
                \move h:.8 v:0
                \beginsegment
                 \unitscale 1
                 \segmentscale .8  % actual scale is 1*2*.7*.8=1.12
                 \fillbox .9
                \endsegment
              \endsegment
             \etg}
$$

This example should be studied carefully. The actual commands that produced
it are 
\beginttverbatim \ninepoint \tt
\1centergraph{\1btg \1unitscale .3 % changes units 
               \1segmentscale 2 
                           % doubles scale of all enclosed segments
                           % actual scale is .3*2=.6
               \1fillbox .5
               \1move h:.8 v:0  % moves to right beyond box
               \1beginsegment
                \1unitscale .5 % now unit scale is .5
                \1segmentscale .7 
                           % reduces scale by .5 -- actual scale is
                           % .5*2*.7=.7 
                \1fillbox .7
                \1move h:.8 v:0
                \1beginsegment
                 \1unitscale 1
                 \1segmentscale .8  % actual scale is 1*2*.7*.8=1.12
                 \1fillbox .9
                \1endsegment
                  % actual scale is back to .5*2*.7=.7
              \1endsegment
                  % actual scale is back to .3*2=.6
             \1etg}
\endttverbatim
Note how the segment scale accumulates while the unitscale does not. Note
also that the scale reverts to that of the enclosing segment after an
|\endsegment|.

It is possible to prevent a segment from being affected by the external
segment scales by inserting the command |\absolutescale|. This is used for
symbols that you wish to have the same size no matter the actual scale of the
graph. 

\sshead{segsum}{Segment Command Summaries}
The following is a list of some of the important points concerning segments:
\beginblockmode
\ext\@|\beginsegment|
\nbr
This begins a graphical segment and a group. All locations within a segment
are relative to the origin of the segment. 

\mbr\ext\@|\endsegment|
\nbr
This terminates a segment and the group. The current location upon leaving
a segment is the segment origin. 

\mbr\ext\@|\unitscale <scale>|
\nbr
The {\tt <scale>} multiplies the present segment and all
enclosed segments unless they have their own |\unitscale|. A |\unitscale| may
be changed at any time within a segment. 

\mbr\ext\@|\segmentscale <scale>|
\nbr
The {\tt <scale>} affects the present segment and all
enclosed segments unless that segment is protected by an |\absolutescale|.
Unlike |\unitscale|, the |\segmentscale| accumulates by multiplying the
current scale value.
The actual scale used for interpreting position numbers is the
product of all of the enclosing segmentscales times the presently active
unitscale. 

\mbr\ext\@|\absoluteposition|
\nbr
All coordinates are to be interpreted with respect to the current origin.
This will either be the graph origin or the origin of the current segment.

\mbr\ext\@|\relativeposition|
\nbr
Coordinate positions are interpreted relative to the previous position. If
the previous command was the drawing of an arc or circle, 
the previous position is
the center of the circle. 
\endblockmode



\shead{slides}{Slide Making}
A specialized set of commands along with appropriate templates exist to
create a transparency. All slides are ``horizontal''. |\fullform| slides
take up an entire 8.5 by 11~in. page and are intended for display.
|\textform| slides are reduced for text. A |\comment <text> //| command is
given that is suppressed when the |\fullform| is used but is printed when
|\textform| is used. The text in the slides uses a general scaled font
family in the file |spfont.tex|. This family is basically a scaled set of
|\tenpoint| fonts. An example is given below:



There are two basic forms, the INRS form and a General form. An example of
the |\inrsform| is given by 


\textform 
\formscale .7

\beginmidinsert
\inrsform

\hslide HOW TO MAKE SLIDES WITH \intex//
\hd{Preliminaries}
\bl 
\lb Input the macro package with {\tt \1input inrssl.tex}.
\lc Decide whether {\tt \1fullform} or {\tt \1textform}.
\last Decide whether {\tt \1inrsform} or {\tt \1genform}.
\el
\hd{Making it.}
\bl 
\lb Use {\tt \1hslide <title>// <body> //}.
\lb Use {``//''} as the delimiter. 
\el
//
\endmidinsert

\beginttverbatim
\1input inrssl % brings in the slide macros
\1textform % sets textform and size
\1formscale .7 % further reduces the size of slide by .7
\1beginmidinsert
\1inrsform % sets the inrs form
\1hslide HOW TO MAKE SLIDES WITH \1intex//
\1hd{Preliminaries}
\1bl 
\1lb Input the macro package with {\1tt \1 1input inrssl.tex}.
\1lc Decide whether {\1tt \1 1fullform} or {\1tt \1 1textform}.
\1last Decide whether {\1tt \1 1inrsform} or {\1tt \1 1genform}.
\1el
\1hd{Making it.}
\1bl 
\1lb Use {\1tt \1 1hslide <title>// <body> //}.
\1lb Use {``//''} as the delimiter. 
\1el
//
\1endmidinsert
\endttverbatim

\beginmidinsert
\topdata = {General Engineering Inc.}
\botdata = {An Important Presentation}
\slidedate = {Jan. 1, 2000}
\genform % this calls the general form that uses these token registers

\hslide This is the title \cr Split here //
\hd{This is an internal header}
\bl 
\lb A listitem  starting with a bullet.
\lc A listitem starting with a circle. 
\last A listitem starting with an asterisk. 
\el
//  % the // terminates a slide

\comment 
It is possible to add a comment to the bottom of the slide which
only shows up when the slide is printed in textform. A comment is
terminated with a {//} //
\endmidinsert

This second slide was given by 
\beginttverbatim
\1topdata = {General Engineering Inc.}
\1botdata = {An Important Presentation}
\1slidedate = {Jan. 1, 2000}
\1genform % the general form that uses these token strings

\1hslide This is the title \1cr Split here //
\1hd{This is an internal header}
\1bl 
\1lb A listitem  starting with a bullet.
\1lc A listitem starting with a circle. 
\1last A listitem starting with an asterisk. 
\1el
//  % the // terminates a slide

\1comment 
It is possible to add a comment to the bottom of the slide which
only shows up when the slide is printed in textform. A comment is
terminated with a {//} //
\endttverbatim


Arbitrary graphics may be included in a slide using normal \TeX Graph
commands. However, the included graphics may not scale correctly when
moving from |\fullform| to |\textform|. Under these conditions is is
necessary to include the test |\iftextform <....> \else <...> \fi| to
switch between the modes. 


The basic commands in the slide package

\beginblockmode
\ext\@|\input inrssl|
\nbr
This command will bring in the INRS\TeX\ slide macro package. 

\mbr\ext\@|\fullform|
\nbr
This chooses the full size form intended for creating the transparency. Any
|\comment| is suppressed.

\mbr\ext\@|\textform|
\nbr
This chooses the text size of the slide and does not suppress the
|\comment| text.

\mbr\ext\@|\formscale <scale>|
\nbr
The entire slide, text, internal graphics, and border, is  scaled 
by the {\tt <scale>} factor. This is done by directly by PostScript. 

\mbr\ext\@|\hslide <title> // <text> //|
This is the basic command for making a slide. The title is terminated with
a // and the slide body is also terminated with a //. 

\mbr\ext\@|\comment <text> //|
This is a comment text which is printed in |\textform| but suppressed in
|\fullform|. 

\mbr\ext\@|\inrsform|
This will give the standard \INRS\ format with its logo on the top and
nothing on the bottom. Text on the bottom may be added using |\botdata| and
|\slidedate|.

\mbr\ext\@|\genform|
\nbr
This is a completely general form where only the border is set.

\mbr\ext\@|\topdata ={<text>}|
This is printed in the |\borderfont| which is the same as the current slide
body font. It is printed on the top line above the border with a current
position at the left hand corner. 

\mbr\ext\@|\botdata = {<b-text>}|
\mbr\ext\@|\slidedate = {<s-date>}|
\nbr 
These two token strings are printed on the bottom line centered and
separated by a --, ie as |<b-text> - <s-date>|. 

\mbr\ext\@|\iftextform <...> \else <...> \fi|
\nbr
This is a switch that allows for different actions between fullform and
textform. |\textform| is .578=1/1.728 the size of |\fullform|.

\mbr\ext\@|\lb \lc \last|
\nbr
These are three short forms for use in listitems starting with $\bullet$,
$\circ$, and {\tt *} respectively. 

\mbr\ext\@|\hd{<text>}|
\nbr
This gives a ``header'' inside a slide. It is equivalent to a |\dssshead|
in normal text. 
\endblockmode


\shead{grcomform}{Some Additional Command Forms -- \TeX Graph}
\beginblockmode
\ext\@|\begintexgraph \btg|
\nbr
This begins \TeX Graph. It will bring in the \TeX Graph macros the first
time it is called. It starts a group and a |\vbox|. 

\mbr\ext\@|\endtexgraph \etg|
\nbr 
Finishes \TeX Graph. It terminates the group started by |\btg|.

\mbr\ext\@|\graphdim <units>|
\nbr 
This sets the default units of \TeX Graph. Any valid \TeX\ |<unit>| is
valid such as {\tt in, cm, pt,} or {\tt bp}. The current default is {\tt
in}.

\mbr\ext\@|\hgraphsize|
\nbr
This maximum horizontal excursion of the graph as computed by \TeX Graph.
It is a dimension. 

\mbr\ext\@|\hgraphskip|
\nbr
This actually |\hskip \hgraphsize|. 

\mbr\ext\@|\hpix <number>/<unit>|
\nbr
This is the horizontal resolution in pixels/unit. The default is 
{\tt 300/in}. 


\mbr\ext\@|\vgraphsize|
\nbr
This maximum vertical excursion of the graph as computed by \TeX Graph.
It is a dimension. 

\mbr\ext\@|\vgraphskip|
\nbr
This actually |\vskip \vgraphsize|. 


\mbr\ext\@|\vpix <number>/<unit>|
\nbr
This is the vertical resolution in pixels/unit. The default is 
{\tt 300/in}. 

\endblockmode
\ejectpage

\done
